\documentclass{article}

\usepackage{amsmath,amssymb}
\usepackage{xurl}
\usepackage[hidelinks]{hyperref}

% Shared commands for notes in docs/paper-gaps/.

\DeclareUnicodeCharacter{2080}{\ifmmode{}_{0}\else\textsubscript{0}\fi}
\DeclareUnicodeCharacter{2081}{\ifmmode{}_{1}\else\textsubscript{1}\fi}
\DeclareUnicodeCharacter{2082}{\ifmmode{}_{2}\else\textsubscript{2}\fi}
\DeclareUnicodeCharacter{2099}{\ifmmode{}_{n}\else\textsubscript{n}\fi}

\newcommand{\Lean}{\textsc{Lean}}
\ExplSyntaxOn
\NewDocumentCommand{\leanid}{m}{
  \ifmmode
    \textsf{\detokenize{#1}}
  \else
    \group_begin:
    \urlstyle{sf}
    \tl_set:Nn \l_tmpa_tl {#1}
    \tl_replace_all:Nnn \l_tmpa_tl {\_} {_}
    \exp_args:NV \path \l_tmpa_tl
    \group_end:
  \fi
}
\ExplSyntaxOff
\providecommand{\tr}{\operatorname{tr}}
\newcommand{\tnleanIssueBase}{https://github.com/LionSR/TNLean/issues/}
\newcommand{\tnleanPrBase}{https://github.com/LionSR/TNLean/pull/}
\newcommand{\tnleanDocBase}{https://sirui-lu.com/TNLean/docs/}
\newcommand{\ghissue}[1]{\href{\tnleanIssueBase#1}{GitHub issue \##1}}
\newcommand{\ghpr}[1]{\href{\tnleanPrBase#1}{PR \##1}}
\newcommand{\doclink}[1]{\href{\tnleanDocBase#1.html}{\path{#1}}}

% Dirac notation and the identity operator, as in blueprint/src/macros/common.tex.
% \providecommand keeps note-local definitions authoritative.
\usepackage{dsfont}
\providecommand{\ket}[1]{|#1\rangle}
\providecommand{\bra}[1]{\langle #1|}
\providecommand{\braket}[2]{\langle #1 | #2 \rangle}
\providecommand{\ketbra}[2]{|#1\rangle\!\langle #2|}
\providecommand{\Id}{\mathds{1}}
\providecommand{\one}{\mathds{1}}
\providecommand{\C}{\mathbb{C}}
\providecommand{\MN}[1]{M_{#1}(\C)}
\providecommand{\MD}{M_D(\C)}

% Verdict marker: \gapnote{<kind>}{<status>}. Kind is the policy
% classification (clarification, local-correction, scope-restriction,
% unfaithful, false-source, open-gap); status is open, wip (elimination
% underway), resolved, or historical. Machine-read by the site index;
% prints a verdict line.
\newcommand{\gapnote}[2]{\par\noindent\textbf{Verdict:} #1 (#2).\par}


\title{Klein Three-Cocycles: Circle and Nonzero-Complex Coefficients}
\author{}
\date{2026-10-01}

\begin{document}
\maketitle
\gapnote{scope-restriction}{open}

\section{The source statement}

Ref.~\cite{GarreRubio2023MPOSym}, subsection ``MPO algebras representing
$\mathbb{Z}_2\times\mathbb{Z}_2$'', states
\begin{align}
    H^3(\mathbb{Z}_2\times\mathbb{Z}_2,U(1)) & \cong\mathbb{Z}_2^3,
    \label{eq:glm23kleinh3_source}
\end{align}
with eight explicit representatives, and that the three diagonal values
$\omega(g,g,g)$, $g\neq e$, of a normalized cocycle distinguish the classes.
Here $U(1)$ carries the trivial action.

\section{What is formalized}

The formal classification uses the multiplicative group $\mathbb{C}^\times$
with the trivial action. Every three-cocycle
$\omega:(\mathbb{Z}_2\times\mathbb{Z}_2)^3\to\mathbb{C}^\times$ is cohomologous,
through a $\mathbb{C}^\times$-valued two-cochain, to exactly one of the eight
representatives, and
\begin{align}
    H^3(\mathbb{Z}_2\times\mathbb{Z}_2,\mathbb{C}^\times) & \cong\mathbb{Z}_2^3.
    \label{eq:glm23kleinh3_complex}
\end{align}
Two cocycles are cohomologous exactly when their three order-two cyclic
invariants $\omega(g,e,g)\,\omega(g,g,g)$ agree.\footnote{
    \leanid{TNLean.Algebra.ScalarThreeCochain.existsUnique\_cohomologousTo\_kleinCocycleFamily},
    \leanid{TNLean.Algebra.ScalarThreeCochain.cohomologousTo\_iff\_klein\_three\_cyclicInvariants},
    and \leanid{TNLean.Algebra.ScalarThreeCochain.kleinH3Equiv} in
    \doclink{TNLean/Algebra/KleinCocycleCompleteness}.}

\section{Missing piece and elimination plan}

The comparison of~\eqref{eq:glm23kleinh3_source}
with~\eqref{eq:glm23kleinh3_complex} is not formalized in degree three; the
existing circle comparison covers degree two. Two routes close it.
First, the proof restricts to circle values: for a $U(1)$-valued cocycle the
normalizing two-cochain $\beta(g,h)=\omega(g,e,e)/\omega(e,e,h)$, the square
root used in the six-entry gauge, and the resulting two-cochain all take
values in $U(1)$, and the eight representatives are $\pm1$-valued. The cyclic
invariants are unchanged by any $\mathbb{C}^\times$-valued coboundary, so they
also separate the $U(1)$ classes. Second, the general comparison
$H^3(G,U(1))\cong H^3(G,\mathbb{C}^\times)$ for finite $G$ follows from
$\mathbb{C}^\times\cong U(1)\times\mathbb{R}_{>0}$, since the uniquely
divisible group $\mathbb{R}_{>0}$ has vanishing positive-degree cohomology
for finite groups.

\textbf{Verdict: scope restriction.} The classification is formalized with
$\mathbb{C}^\times$ coefficients; the circle-valued statement is not.

\bibliographystyle{alpha}
\bibliography{references}

\end{document}
